% --------------------------------- PARTIE1----------
\block{\textsc{Introduction}}{
    In many supply chains, demand appears stable at the consumer level — yet grows increasingly volatile further upstream. This poster investigates why, and how machine learning can fix it.

    \vspace{0.4cm}

    \begin{itemize}
        \item \textbf{Multi-echelon supply chain:} consumer $\to$ retailer $\to$ wholesaler $\to$ distributor $\to$ producer (steel products).
        \item \textbf{Downstream demand:} stable and forecastable close to the consumer.
        \item \textbf{Upstream orders:} amplified volatility — peaks and troughs disconnected from real demand.
        \item \textbf{Consequences:} extra costs (12.5–25\%), overcapacity, excess or short stocks.
        \item \textbf{This phenomenon:} the \textbf{bullwhip effect}.
    \end{itemize}
    \begin{center}        
         \includegraphics[width=25cm]{images/Bullwhip_Effect.png}
    \end{center}
    \vspace{0.2cm}

    $\rightarrow$ This poster shows that hybrid ML models (ARIMAX + Neural Networks + exogenous indicators) significantly outperform traditional forecasting methods, improving forecast accuracy, inventory turns, and cash-conversion cycle.

}
      % --------------------------------- PARTIE 2  ----------    
\block{\textsc{Bullwhip effect}}{
    Here is the formula which describes the "Bullwhip effect"
    \vspace{1cm}
    \[
        D_t = d + \rho D_{t-1} + u_t
    \]
    \vspace{1cm}
    \begin{itemize}
        \setlength{\leftskip}{1.2cm}
        \item $\mathbf{D_t}$: the demand observed today, for example this week's sales.
        \item $\mathbf{D_{t-1}}$: the demand from the previous period, for example last week's sales.
        \item $\mathbf{d}$: the baseline level of demand, that is, the ``normal'' demand around which fluctuations occur.
        \item $\mathbf{\rho}$: the memory effect of demand; if $\mathbf{\rho > 0}$, a recent increase in demand tends to be followed by another increase, at least partially.
        \item $\mathbf{u_t}$: the random noise term, that is, everything that makes demand move in an unexpected way.
    \end{itemize}
    \vspace{2cm}
    \begin{quote}
    Today's demand depends on an average level, a part of past demand, and an unpredictable shock.
    \end{quote}
} 